% Ad Familiares 14.1 (ad-familiares-14-01)
% Source: https://raw.githubusercontent.com/PerseusDL/canonical-latinLit/master/data/phi0474/phi056/phi0474.phi056.perseus-lat1.xml
% Fetched by scripts/fetch_latin.py

% Dateline (Perseus): Scr. Dyrrhachi a. d. vi K. Dec. a. 696 (48).

\ciceroLetterOpener{TVLLIVS TERENTIAE SVAE, TVLLIOLAE SVAE, CICERONI SVO S. D.}

\ciceroSection{1} et litteris multorum et sermone omnium perfertur ad me incredibilem tuam virtutem et fortitudinem esse teque nec animi neque corporis laboribus defetigari. me miserum! te ista virtute, fide, probitate, humanitate in tantas aerumnas propter me incidisse, Tulliolamque nostram, ex quo patre tantas voluptates capiebat, ex eo tantos percipere luctus! nam quid ego de Cicerone dicam? qui cum primum sapere coepit, acerbissimos dolores miseriasque percepit. quae si, tu ut scribis, 'fato facta' putarem, ferrem paulo facilius ; sed omnia sunt mea culpa commissa, qui ab iis me amari putabam qui invidebant, eos non sequebar qui petebant.

\ciceroSection{2} quod si nostris consiliis usi essemus neque apud nos, tantum valuisset sermo aut stultorum amicorum aut improborum, beatissimi viveremus. nunc quoniam sperare nos amici iubent, dabo operam ne mea valetudo tuo labori desit. res quanta sit intellego quantoque fuerit facilius manere domi quam redire ; sed tamen si omnis tr. pl. habemus, si Lentulum tam studiosum quam videtur, si vero etiam Pompeium et Caesarem, non est desperandum.

\ciceroSection{3} de familia quo modo placuisse scribis amicis faciemus. de loco nunc quidem iam abiit pestilentia, sed quam diu fuit me non attigit. Plancius, homo officiosissimus, me cupit esse secum et adhuc retinet. ego volebam loco magis deserto esse in Epiro, quo neque Piso veniret nec milites, sed adhuc Plancius me retinet ; sperat posse fieri ut mecum in Italiam decedat. quem ego diem si videro et si in vestrum complexum venero ac si et vos et me ipsum reciperaro, satis magnum mihi fructum videbor percepisse et vestrae pietatis et meae.

\ciceroSection{4} Pisonis humanitas, virtus, amor in omnis nos tantus est, ut nihil supra possit. utinam ea res ei voluptati sit ! gloriae quidem video fore. de Quinto fratre nihil ego te accusavi, sed vos, cum praesertim tam pauci sitis, volui esse quam coniunctissimos.

\ciceroSection{5} quibus me voluisti agere gratias egi et me a te certiorem factum esse scripsi. quod ad me, mea Terentia, scribis te vicum vendituram, quid, obsecro te, me miserum! quid futurum est? et si nos premet eadem fortuna, quid puero misero fiet? non queo reliqua scribere ; tanta vis lacrimarum est ; neque te in eundem fletum adducam ; tantum scribo : si erunt in officio amici, pecunia non deerit ; si non erunt, tu efficere tua pecunia non poteris. per fortunas miseras nostras, vide ne puerum perditum perdamus ; cui si aliquid erit ne egeat, mediocri virtute opus est et mediocri fortuna ut cetera consequatur.

\ciceroSection{6} fac valeas et ad me tabellarios mittas, ut sciam quid agatur et vos quid agatis. mihi omnino iam brevis exspectatio est, Tulliolae et Ciceroni salutem dic. valete . D. a. d. vi K. Decembr. Dyrrhachi.

\ciceroSection{7} Dyrrhachium veni, quod et libera civitas est et in me officiosa et proxima Italiae ; sed si offendet me loci celebritas, alio me conferam, ad te scribam.
